\subsection{The relative pro-\texorpdfstring{\(p\)}{p} fibre over an infinite coprime base}

\begin{statement}{D.43}{Proved}{Compatible relative lifts over an actual pro-prime-to-$p$ quotient}
Let \(q:\widehat F\twoheadrightarrow D\) be an actual
profinite quotient, let \(K=\ker q\), and suppose that
every finite quotient of \(D\) has order prime to \(p\).
Fix \(\alpha\in\operatorname{Aut}(D)\) and actual
elements \(h_i\in D\) such that
\[
                  \alpha(d_i)=h_i d_i h_i^{-1}.
 \tag{D.43.1}
\]
For a profinite group \(H\), write
\[
 O^p(H)=
 \bigcap_{\substack{N\triangleleft_{\mathrm{open}}H\\
                         H/N\ {\rm a\ finite}\ p{\rm\text{-}group}}}N.
 \tag{D.43.2}
\]
Then the actual source extension
\[
 E_p(D)=\widehat F/O^p(K),\qquad
        1\longrightarrow K^{(p)}
          \longrightarrow E_p(D)\longrightarrow D\longrightarrow1
 \tag{D.43.3}
\]
has a multiplier-one peripheral automorphism above
\(\alpha\), with peripheral conjugator projections
exactly \(h_i\).

Choose any cofinal decreasing sequence of
characteristic open normal subgroups \(W_j\) of \(D\),
put \(D_j=D/W_j\), \(K_j=q^{-1}(W_j)\), and
\(E_{j,p}=\widehat F/O^p(K_j)\).
The full literal fibre in \textup{(D.43.3)} is
canonically the inverse limit of the complete
finite-base fibres of D.42:
\[
 \begin{split}
 E_p(D)&\ \simeq\ \varprojlim_j E_{j,p},\\
 \mathscr X(D,\alpha,h)&\ \simeq\
             \varprojlim_j\mathscr X(D_j,\alpha_j,h_j).
 \end{split}
 \tag{D.43.4}
\]
Both comparisons retain the original map from
\(\widehat F\).
The fibres are compact and the inverse limit
is nonempty.
Every generating source map
\(\widehat F\to E_p(D)\) above \(\alpha q\)
factors through an automorphism of this same
\(E_p(D)\); its three peripheral conditions
are exactly those in the displayed fibre.

In particular let \(K_n=\ker(\widehat F\to P_n)\)
for the D.37 quotient.
For every odd prime \(p\ne5\), every D.37
parameter has a multiplier-one peripheral
realization on the single actual quotient
\[
       \widehat F/
             O^p\bigl(O^2(K_n)\bigr).
 \tag{D.43.5}
\]
It induces a chosen full relative pro-two
lift from D.41, and hence induces the given
\(\beta_Q\) on \(P_n\).
Its maximal pro-five quotient is exactly
\(P_n\).
The resulting automorphisms are compatible
through all the finite two-primary stages
and their relative \(p\)-central stages.
\end{statement}

\begin{proof}
We first prove the residual and quotient
comparisons without a finiteness assumption
on \(K\).

If \(H\) is a closed subgroup of \(H'\), then
\[
                         O^p(H)\subset O^p(H').
 \tag{D.43.6}
\]
Indeed every finite \(p\)-quotient map of
\(H'\), restricted to \(H\), kills \(O^p(H)\).
If \(f:H\twoheadrightarrow J\) is a continuous
epimorphism, then
\[
                         f(O^p(H))=O^p(J).
 \tag{D.43.7}
\]
One inclusion follows by pulling back every
finite \(p\)-quotient of \(J\).
For the other, the quotient
\(J/f(O^p(H))\) is a continuous quotient of
the pro-\(p\) group \(H/O^p(H)\).
The subgroup \(f(O^p(H))\) is closed because
it is a compact image, so this is an actual
pro-\(p\) quotient.
It must contain the defining residual
\(O^p(J)\) in its kernel.
This proves \textup{(D.43.7)} with its exact
maps.

The group \(D\) is finitely generated,
as a quotient of \(\widehat F\).
It has only finitely many subgroups of each
given finite index: each such subgroup
defines a permutation action, and the
images of a fixed finite generating set
give only finitely many such actions.
Intersections of all open normal subgroups
of index at most \(j\) therefore give a
cofinal decreasing characteristic sequence
\(W_j\).
In particular \(\alpha\) induces the
compatible automorphisms \(\alpha_j\)
of \(D_j\).
By the hypothesis on \(D\), every \(D_j\)
has order prime to \(p\).

\paragraph{The residual equality and the first comparison.}
The inclusions \textup{(D.43.6)} give
\[
 O^p(K)\subset O^p(K_k)\subset O^p(K_j)
                         \quad(k\geq j).
 \tag{D.43.8}
\]
Every displayed subgroup is normal in
\(\widehat F\): its ambient \(K_j\), or
\(K\), is normal, and its residual is
characteristic in that ambient group.
We claim the exact equality
\[
                   \bigcap_j O^p(K_j)=O^p(K).
 \tag{D.43.9}
\]
To prove the remaining inclusion, take an
arbitrary finite quotient \(G\) of
\(\widehat F/O^p(K)\), and let
\(f:\widehat F\twoheadrightarrow G\)
be its source map.
The image \(N=f(K)\) is a finite
\(p\)-group, because the restriction to
\(K\) factors through \(K^{(p)}\).
There is an actual induced quotient map
\[
                          D\longrightarrow G/N.
 \tag{D.43.10}
\]
Its kernel is open, so some \(W_j\)
is contained in that kernel.
It follows that \(f(K_j)\subset N\).
The restriction of \(f\) to \(K_j\)
then has \(p\)-group image and kills
\(O^p(K_j)\).
Consequently the left-hand side of
\textup{(D.43.9)} is killed in every
finite quotient of
\(\widehat F/O^p(K)\), which proves
the reverse inclusion.
Notice that this residual equality
itself holds for any closed normal \(K\)
and these open approximations; the
prime-to-\(p\) assumption is used for
the nonempty finite-base fibres.

The source quotient maps now give
\[
 \widehat F/O^p(K)\longrightarrow
                  \varprojlim_j\widehat F/O^p(K_j).
 \tag{D.43.11}
\]
Its kernel is trivial by
\textup{(D.43.9)}.
It is onto: each compatible tuple of
cosets has nested nonempty closed
inverse-image sets in compact
\(\widehat F\).
Their intersection provides a
simultaneous source lift.
Thus \textup{(D.43.11)} is an
isomorphism of profinite groups,
not merely an identification of
their finite images.

\paragraph{Every upper fibre has an actual descent map.}
For \(k\geq j\), let
\(\pi_{k,j}:E_{k,p}\to E_{j,p}\)
be the map supplied by
\textup{(D.43.8)}.
Inside \(E_{k,p}\) put
\[
 M_{j,k}
   =\ker(E_{k,p}\longrightarrow D_j)
   =K_j/O^p(K_k).
 \tag{D.43.12}
\]
The last equality uses the actual
image of \(K_j\) under the source map.
Applying \textup{(D.43.7)} to
\(K_j\twoheadrightarrow M_{j,k}\)
gives
\[
 O^p(M_{j,k})
       =O^p(K_j)/O^p(K_k)
       =\ker\pi_{k,j}.
 \tag{D.43.13}
\]
An automorphism of \(E_{k,p}\)
above \(\alpha_k\) preserves
\(M_{j,k}\), because
\(\alpha_k\) preserves the kernel
of \(D_k\to D_j\).
It therefore preserves the
characteristic subgroup
\(O^p(M_{j,k})\).
Formula \textup{(D.43.13)} proves
that it descends to an actual
automorphism of \(E_{j,p}\).
The descended map is above
\(\alpha_j\); every one of its
peripheral conjugators is the
image of an upper conjugator.
It has projection \(h_{i,j}\)
when the upper one has projection
\(h_{i,k}\).

Here and below the fibre is the
literal compact set
\[
 \begin{split}
 \mathscr X_j=\{\sigma\in\operatorname{Aut}(E_{j,p}):&
 \ \bar q_j\sigma=\alpha_j\bar q_j,\ \\
 &\sigma(X_i)=a_iX_i a_i^{-1}
 \text{ for some }a_i\in E_{j,p}
 \text{ with }\bar q_j(a_i)=h_{i,j},
 \ 1\leq i\leq3\}.
 \end{split}
 \tag{D.43.14}
\]
The third source element is
\(X_3=(X_1X_2)^{-1}\), as throughout
the paper.
The preceding descent is a
continuous map
\(\mathscr X_k\to\mathscr X_j\).
The maps compose exactly because
they are induced by the same
source quotient maps.
Thus \textup{(D.43.14)} is an
actual inverse system.

\paragraph{Compactness, nonemptiness, and the three exact conjugators.}
Every \(E_{j,p}\) is finitely
generated, being a quotient of
\(\widehat F\).
For completeness its automorphism
group is compact.
Take a characteristic cofinal
sequence of finite quotients,
formed by bounded-index
intersections as above.
An automorphism gives compatible
automorphisms of those finite
quotients.
A compatible tuple, with its
compatible inverses, defines
an automorphism of the original
profinite group.
The compatibility and subgroup
preservation conditions are
closed in the product of the
finite automorphism sets.
This identifies the automorphism
group with a closed profinite
subset and proves compactness.

For each \(i\), the simultaneous
condition on \(\sigma,a_i\) in
\textup{(D.43.14)} is closed.
Projecting its compact solution
set to the automorphism factor
shows that the existential
conjugator condition is also
closed.
Hence \(\mathscr X_j\) is
compact, and D.42 proves it
nonempty.
Every finite collection of
inverse-limit equations has
a solution: choose an element
of \(\mathscr X_k\) at an
index at least all the indices
in that collection and use
its actual descents.
Compactness of the product
therefore gives
\[
                         \varprojlim_j\mathscr X_j
                                  \ne\varnothing .
 \tag{D.43.15}
\]
No surjectivity of the
transition maps is needed.

A compatible tuple
\((\sigma_j)_j\) defines an
automorphism \(\sigma\) of
\(E_p(D)\) through
\textup{(D.43.11)}.
Its inverse is the tuple of
the compatible \(\sigma_j^{-1}\).
It lies above \(\alpha\).
To recover a conjugator with
exact projection \(h_i\),
consider in \(E_p(D)\) the
closed sets
\[
 \mathcal A_{i,j}=
 \{a:\ q(a)\equiv h_i\pmod {W_j},\
      \sigma_j(X_i)=\pi_j(a)X_i\pi_j(a)^{-1}\}.
 \tag{D.43.16}
\]
Each is nonempty: lift a
conjugator from
\textup{(D.43.14)} through
the surjection \(\pi_j\).
The sets are decreasing
by the descent equations.
Their intersection is
nonempty by compactness.
Any element \(a_i\) in it
satisfies \(q(a_i)=h_i\)
and
\(\sigma(X_i)=a_iX_i a_i^{-1}\),
since all finite-base
pro-\(p\) projections detect
equality in
\textup{(D.43.11)}.
This establishes the three
literal peripheral conditions.

Conversely an automorphism
of \(E_p(D)\) above \(\alpha\)
preserves each inverse image
of \(W_j\).
Applying
\textup{(D.43.7)} to
\(K_j\to K_j/O^p(K)\)
identifies the pro-\(p\)
residual of this inverse
image with
\(\ker(E_p(D)\to E_{j,p})\).
Thus it descends to every
stage, with the specified
peripheral conditions.
This is inverse to the
construction just given.
The two maps are continuous
in their compact
automorphism topologies.
They prove the second
isomorphism in
\textup{(D.43.4)} for the
complete fibre.

\paragraph{All source maps and the fixed five-group quotient.}
Let \(f:\widehat F\twoheadrightarrow E_p(D)\)
satisfy \(q f=\alpha q\).
It sends \(K\) into the
normal pro-\(p\) kernel
\(K/O^p(K)\).
Hence it kills \(O^p(K)\)
and gives a surjective
endomorphism of the same
\(E_p(D)\).
Every finitely generated
profinite group is Hopfian
by the following finite
argument.
The inverse-image map of
a surjective endomorphism
permutes the finite set of
open normal subgroups of
each bounded index:
it is injective on that
set and preserves indices.
It therefore preserves
their intersection and
induces an automorphism
on each resulting finite
quotient.
Their intersection is
trivial, so the endomorphism
is injective.
Apply this to the
two-generated quotient
\(E_p(D)\).
There is no finite-rank
assumption on \(K^{(p)}\)
in this argument.
This proves the
all-source factorization
and its peripheral
comparison.

Finally set
\[
                 D=\widehat F/O^2(K_n).
 \tag{D.43.17}
\]
Its kernel over \(P_n\)
is \(K_n^{(2)}\), and
every finite quotient
has order \(2^a5^b\).
For odd \(p\ne5\) it is
therefore a pro-prime-to-\(p\)
base in the precise sense
of the statement.
Choose the actual
multiplier-one automorphism
\(\alpha\) of this \(D\)
supplied by D.41 above
\(\beta_Q\), and choose
its actual three
conjugators \(h_i\in D\).
The kernel of
\(\widehat F\to D\) is
exactly \(O^2(K_n)\).
Substitution into
\textup{(D.43.3)} gives
the quotient
\textup{(D.43.5)} with
its literal source map.
The finite lower-two-central
quotients \(D_{n,j}\) of
D.42 are cofinal here:
their pro-two kernel
series terminates in
every finite two-group
image.
They are preserved by
the chosen \(\alpha\).
The proof above only
requires that stability,
so it yields precisely
the compatible lift
through these stages
and all their finite
relative \(p\)-central
quotients.

Every pro-five quotient
of \textup{(D.43.5)}
kills first its normal
pro-\(p\) kernel and
then the normal pro-two
kernel of
\textup{(D.43.17)}.
It factors through \(P_n\),
which remains an actual
quotient.
This proves the exact
maximal pro-five
identification and
the retained action
\(\beta_Q\).
\end{proof}

D.43 supplies the
across-stage compatibility
left open within D.42.
It concerns a complete
relative pro-\(p\) kernel
over a base all of whose
finite quotients have
order prime to \(p\).
A further step in which
the same prime already
occurs in the base
requires additional
source comparisons.
The full D.1--D.2
leading-image question
continues to include
those finite quotients.
